\documentclass[a4paper,12pt]{article}
\usepackage[T2A]{fontenc}
\usepackage[utf8]{inputenc}
\usepackage[russian]{babel}
\usepackage{amsmath,amssymb,mathtools}
\usepackage{helvet}
\renewcommand{\familydefault}{\sfdefault}
\usepackage{icomma}
\usepackage{geometry}
\geometry{left=24mm,right=22mm,top=24mm,bottom=24mm}
\usepackage{microtype}
\usepackage{tikz}
\usetikzlibrary{calc,arrows.meta,positioning,shapes.geometric}
\usepackage{tabularx,booktabs,array}
\usepackage{enumitem}
\usepackage{parskip}
\usepackage{fancyhdr}
\usepackage{setspace}
\usepackage{xcolor}
\usepackage{hyperref}
\hypersetup{hidelinks,pdfauthor={Лёвин Артём Александрович},pdftitle={Чек-лист. Финансовая математика: вклады и сложные проценты}}
\definecolor{accent}{HTML}{008080}
\definecolor{darkaccent}{HTML}{004D4D}
\definecolor{textgray}{HTML}{333333}
\color{textgray}
\onehalfspacing
\setlength{\headheight}{15pt}
\pagestyle{fancy}
\fancyhf{}
\fancyhead[L]{\small Вклады и сложные проценты}
\fancyhead[R]{\small Ксения Клыкова}
\fancyfoot[C]{\thepage}
\renewcommand{\headrulewidth}{0.4pt}
\renewcommand{\footrulewidth}{0pt}
\setlist[itemize]{leftmargin=1.65em,itemsep=1pt,topsep=3pt}
\setlist[enumerate]{leftmargin=1.8em,itemsep=5pt,topsep=5pt}
\newcommand{\lessonsection}[1]{\section{#1}\vspace{-0.2em}\noindent\textcolor{accent}{\rule{\linewidth}{0.55pt}}\vspace{0.2em}}
\newcommand{\keyidea}[1]{\par\noindent\textcolor{darkaccent}{\textbf{Главная идея.}} #1\par}
\newcommand{\exampletitle}[1]{\par\smallskip\noindent\textcolor{darkaccent}{\textbf{#1}}\par}
\begin{document}
\thispagestyle{empty}
\begin{titlepage}
\centering
\vspace*{25mm}
{\Large\color{darkaccent} МАТЕМАТИКА\par}
\vspace{14mm}
{\huge Чек-лист\par}
\vspace{10mm}
{\LARGE\color{accent} Финансовая математика\\[3mm] Вклады и сложные проценты\par}
\vspace{14mm}
{\large Подготовка к ЕГЭ по профильной математике\par}
\vfill
{\Large Ксения Клыкова\par}
\vspace{10mm}
{\large 26 октября 2026 г.\par}
\vspace{11mm}
{\normalsize Лёвин Артём Александрович\\ эксклюзивно для Ксении Клыковой\par}
\vspace{12mm}
\begin{tikzpicture}
\draw[accent,line width=0.9pt] (-5.4,0) .. controls (-4.0,0.85) and (-2.1,-0.7) .. (0,0)
.. controls (1.9,0.75) and (3.7,-0.7) .. (5.4,0.05);
\fill[darkaccent] (0,0) circle (1.3pt);
\end{tikzpicture}
\vspace{9mm}
\end{titlepage}
\setcounter{page}{1}
\lessonsection{Что происходит с деньгами на вкладе}
Банк начисляет проценты на сумму, которая находится на вкладе к моменту начисления. Если начисленные проценты остаются на счёте, то на следующий год они сами приносят доход. Это \textbf{капитализация процентов}.

Пусть начальная сумма равна $S$, ставка составляет $r\%$ годовых, а число лет равно $n$. Удобно перейти к двум коэффициентам:
\[
 p=\frac{r}{100},\qquad k=1+p=1+\frac{r}{100}.
\]
Здесь $p$ --- ставка в долях единицы, $k$ --- \textbf{повышающий коэффициент}. Например, при ставке $10\%$ получаем $p=0{,}1$ и $k=1{,}1$.

\exampletitle{Пример 1. Вклад 100 рублей под 10\% годовых на три года}
\begin{center}
\renewcommand{\arraystretch}{1.30}
\begin{tabular}{@{}c r r r@{}}
\toprule
Год & В начале года, руб. & Проценты, руб. & В конце года, руб.\\
\midrule
1 & $100$ & $10$ & $110$\\
2 & $110$ & $11$ & $121$\\
3 & $121$ & $12{,}1$ & $133{,}1$\\
\bottomrule
\end{tabular}
\end{center}
\keyidea{Проценты второго года вычисляются от $110$ рублей, третьего --- от $121$ рубля. Каждый год вся накопленная сумма умножается на $k$.}

\lessonsection{Вывод формулы сложных процентов}
Выполним ту же операцию с буквами, как на занятии:
\begin{center}
\renewcommand{\arraystretch}{1.30}
\begin{tabular}{@{}c c c c@{}}
\toprule
Год & Начальная сумма & Начислено процентов & Конечная сумма\\
\midrule
1 & $S$ & $Sp$ & $S+Sp=Sk$\\
2 & $Sk$ & $Skp$ & $Sk(1+p)=Sk^2$\\
3 & $Sk^2$ & $Sk^2p$ & $Sk^2(1+p)=Sk^3$\\
\bottomrule
\end{tabular}
\end{center}
Показатель степени совпадает с числом лет. Поэтому
\[
\boxed{\displaystyle S_n=Sk^n=S\left(1+\frac{r}{100}\right)^n.}
\]
Формула применяется, если \textbf{ставка постоянна}, проценты начисляются один раз в конце каждого года и \textbf{нет пополнений и снятий}. Индекс $n$ обозначает количество годовых начислений.

\exampletitle{Пример 2. Сумма 200 рублей, ставка 20\%, три года}
\[
S_3=200\cdot1{,}2^3=200\cdot1{,}728=345{,}6\ \text{руб.}
\]
\textcolor{darkaccent}{\textbf{Проверка смысла:}} при положительной ставке итоговая сумма больше первоначальной.

\lessonsection{Пополнение в начале года}
Если в начале года вносится дополнительная сумма, \textbf{сначала увеличивается остаток}, затем проценты начисляются на новую сумму. Разделять эти действия важно.

Пусть $S$ --- первоначальный вклад, а в начале второго и третьего годов дополнительно вносят по $a$ рублей. Тогда:
\[
\begin{aligned}
\text{конец 1-го года:}\quad&B_1=Sk,\\
\text{конец 2-го года:}\quad&B_2=(Sk+a)k=Sk^2+ak,\\
\text{конец 3-го года:}\quad&B_3=((Sk+a)k+a)k.
\end{aligned}
\]
Раскрываем скобки и получаем основную формулу занятия:
\[
\boxed{\displaystyle B_3=Sk^3+ak^2+ak.}
\]
Первое слагаемое --- рост исходной суммы. Два других отражают доходность дополнительных взносов: взнос второго года успевает пройти \textbf{два} начисления, а взнос третьего --- только \textbf{одно}.

\exampletitle{Пример 3. Вклад 120 млн руб. под 15\% на три года}
В начале второго и третьего годов вносят по $10$ млн руб. Получаем $k=1{,}15$:
\begin{center}
\renewcommand{\arraystretch}{1.23}
\begin{tabular}{@{}c r r r r@{}}
\toprule
Год & Доп. взнос & В начале года & Проценты & В конце года\\
\midrule
1 & $0$ & $120$ & $18$ & $138$\\
2 & $10$ & $148$ & $22{,}2$ & $170{,}2$\\
3 & $10$ & $180{,}2$ & $27{,}03$ & $207{,}23$\\
\bottomrule
\end{tabular}
\end{center}
\noindent\smallskip
Все числа таблицы выражены в \textbf{миллионах рублей}. По формуле:
\[
120\cdot1{,}15^3+10\cdot1{,}15^2+10\cdot1{,}15
=207{,}23.
\]

\lessonsection{Снятие денег и срок каждого взноса}
Если в начале года снимают деньги, вместо добавления соответствующей суммы нужно её \textbf{вычесть}. Для регулярных операций по $a$ в начале каждого года, кроме первого, полезно запомнить три случая:
\[
\begin{array}{ll}
\text{2 года:}& B_2=Sk^2\ \pm ak,\\[2mm]
\text{3 года:}& B_3=Sk^3\ \pm ak^2\ \pm ak,\\[2mm]
\text{4 года:}& B_4=Sk^4\ \pm ak^3\ \pm ak^2\ \pm ak.
\end{array}
\]
Знак «$+$» выбирают для пополнений, знак «$-$» --- для снятий. При каждом снятии сумма на счёте должна оставаться достаточной для операции.

\exampletitle{Пример 4. Четыре года со снятием денег}
Внесено $100$ тыс. руб. под $10\%$. В начале второго, третьего и четвёртого годов снимают по $5$ тыс. руб. Тогда
\[
\begin{aligned}
B_4&=100\cdot1{,}1^4-5(1{,}1^3+1{,}1^2+1{,}1)\\
&=146{,}41-18{,}205=128{,}205\ \text{тыс. руб.}
\end{aligned}
\]
\textbf{Ответ в рублях:} $128\,205$ руб., поскольку $128{,}205\cdot1000=128\,205$.

\lessonsection{Один универсальный порядок расчёта}
Когда суммы пополнений и снятий различаются, удобнее действовать по годам. Обозначим $B_0=S$. Пусть $c_t$ --- операция \emph{в начале} года $t$ (пополнение положительно, снятие отрицательно; обычно $c_1=0$). Тогда
\[
\boxed{\displaystyle B_t=(B_{t-1}+c_t)k,\qquad t=1,2,\ldots,n.}
\]
Например, при $S=200$ тыс. руб., $k=1{,}1$, пополнении на $20$ тыс. руб. в начале второго года и снятии $10$ тыс. руб. в начале третьего:
\[
B_3=200\cdot1{,}1^3+20\cdot1{,}1^2-10\cdot1{,}1
=279{,}4\ \text{тыс. руб.}
\]
\keyidea{Каждая операция приносит доход (или уменьшает будущую сумму) ровно за то количество начислений, которое остаётся после неё, включая конец года операции.}

\lessonsection{Обратные задачи: найти вклад или ставку}
\exampletitle{Пример 5. Найти первоначальную сумму}
Вклад на три года под $20\%$. В начале второго и третьего годов добавляют по $10$ тыс. руб. В конце третьего года накоплено $173{,}28$ тыс. руб. Пусть начальная сумма равна $S$:
\[
\begin{aligned}
173{,}28&=S\cdot1{,}2^3+10\cdot1{,}2^2+10\cdot1{,}2,\\
173{,}28&=1{,}728S+26{,}4,\\
S&=\frac{173{,}28-26{,}4}{1{,}728}=85.
\end{aligned}
\]
\textbf{Ответ:} $85$ тыс. руб.

\exampletitle{Пример 6. Найти годовую процентную ставку}
Вклад $500$ тыс. руб. за три года вырос до $665{,}5$ тыс. руб.; пополнений и снятий нет. Используем повышающий коэффициент:
\[
500k^3=665{,}5 \quad\Longrightarrow\quad
k^3=1{,}331=1{,}1^3 \quad\Longrightarrow\quad k=1{,}1.
\]
Значит, $p=k-1=0{,}1$, а $r=100p=10$. \textbf{Ответ:} $10\%$ годовых.

Для вкладов с неотрицательной ставкой $k\ge1$; при любых ставках, допустимых для данного вклада, должно выполняться $k>0$. Это помогает отсекать бессмысленные значения при вычислениях.

\lessonsection{Типичные ошибки и итоговая памятка}
\begin{itemize}
\item \textbf{Процент и коэффициент.} При $15\%$ используйте $p=0{,}15$ и $k=1{,}15$, а не $k=15$ или $k=0{,}15$.
\item \textbf{Начало или конец года.} Дополнительный взнос, внесённый \emph{в начале} года, уже участвует в ближайшем начислении процентов.
\item \textbf{Лишняя или пропущенная степень.} Взнос в начале второго года трёхлетнего вклада умножается на $k^2$; в начале третьего --- на $k$.
\item \textbf{Снятие.} Снятая сумма уменьшает и текущий остаток, и будущие проценты; перед соответствующим членом формулы стоит минус.
\item \textbf{Единицы.} Сумма в тысячах рублей переводится в рубли умножением на $1000$; в миллионах --- на $1\,000\,000$.
\item \textbf{Обратная задача.} Сначала найдите $k^n$ или коэффициент при $S$. Подставлять длинное выражение с процентами лучше после упрощения.
\end{itemize}
\begin{center}
\renewcommand{\arraystretch}{1.25}
\begin{tabularx}{\linewidth}{@{}p{0.40\linewidth}X@{}}
\toprule
Ситуация & Что делать\\
\midrule
Нет операций со вкладом & $S_n=Sk^n$.\\
Дополнительный взнос & Добавить сумму до начисления процентов.\\
Снятие & Вычесть сумму до начисления процентов.\\
Неизвестна начальная сумма & Выразить $S$ из уравнения.\\
Неизвестна ставка & Найти $k$, затем $r=100(k-1)$.\\
\bottomrule
\end{tabularx}
\end{center}

\clearpage
\thispagestyle{fancy}
\begin{center}
{\Large\color{darkaccent}\textbf{Алгоритм: расчёт вклада по годам}}\par
\vspace{4mm}
\end{center}
\vspace{8mm}
\begin{center}
\begin{tikzpicture}[
>=Stealth,
term/.style={ellipse,draw=accent,line width=0.8pt,align=center,minimum width=55mm,minimum height=12mm,font=\small},
act/.style={rounded corners=2pt,draw=accent,line width=0.75pt,align=center,text width=102mm,minimum height=14mm,font=\small,inner sep=5pt},
dec/.style={diamond,aspect=2.9,draw=darkaccent,line width=0.8pt,align=center,text width=28mm,inner sep=1.5pt,font=\small},
arr/.style={->,line width=0.85pt,draw=darkaccent}
]
\node[term] (start) {Начало: известны $S$, $r$, $n$ и операции $c_t$};
\node[act,below=12mm of start] (init) {Найти $k=1+\dfrac{r}{100}$; положить $B_0=S$, $t=1$};
\node[act,below=13mm of init] (input) {В начале года $t$ выполнить операцию:\\$A_t=B_{t-1}+c_t$};
\node[act,below=13mm of input] (interest) {В конце года начислить проценты:\\$B_t=A_t\,k$};
\node[dec,below=15mm of interest] (test) {$t=n$?};
\node[act,below=15mm of test] (answer) {Записать итоговую сумму $B_n$\\и проверить единицы ответа};
\node[term,below=12mm of answer] (end) {Конец};
\node[act,right=22mm of test,text width=31mm,minimum height=12mm] (inc) {Увеличить $t$ на $1$};
\draw[arr] (start)--(init);
\draw[arr] (init)--(input);
\draw[arr] (input)--(interest);
\draw[arr] (interest)--(test);
\draw[arr] (test)--node[left,font=\small]{да}(answer);
\draw[arr] (answer)--(end);
\draw[arr] (test.east)--node[above,font=\small]{нет}(inc.west);
\draw[arr] (inc.north) |- (input.east);
\end{tikzpicture}
\end{center}
\clearpage
\lessonsection{Тренировочный блок --- Путь А}
\begingroup
\small
\setstretch{1.18}
\setlist[enumerate]{leftmargin=1.8em,itemsep=2pt,topsep=2pt,parsep=0pt}
\textbf{Условия для всех задач:} проценты начисляются один раз в конце каждого года, а названные пополнения и снятия происходят в начале соответствующего года. Ставка постоянна. В заданиях с ответом в рублях обратите внимание на единицы исходной суммы.

\subsection*{Базовый уровень}
\begin{enumerate}
\item В банк положили $1000$ руб. под $10\%$ годовых. Какая сумма будет через один год?
\item Вклад составляет $80$ тыс. руб. под $5\%$ годовых. Найдите сумму через два года без дополнительных операций. Ответ дайте в тысячах рублей.
\item Вклад $200$ тыс. руб. размещён на три года под $20\%$ годовых. Найдите конечную сумму в тысячах рублей.
\end{enumerate}
\subsection*{Средний уровень}
\begin{enumerate}[resume]
\item Вклад $100$ тыс. руб. размещён на три года под $10\%$ годовых. В начале второго и третьего годов его пополняют на $10$ тыс. руб. Найдите конечную сумму в тысячах рублей.
\item Вклад $100$ тыс. руб. размещён на три года под $10\%$ годовых. В начале второго и третьего годов снимают по $5$ тыс. руб. Найдите конечную сумму в тысячах рублей.
\item Вклад $50$ тыс. руб. размещён на четыре года под $20\%$ годовых. В начале второго, третьего и четвёртого годов добавляют по $5$ тыс. руб. Найдите сумму к концу четвёртого года в тысячах рублей.
\end{enumerate}
\subsection*{Уровень выше среднего}
\begin{enumerate}[resume]
\item На три года под $10\%$ годовых открыт вклад неизвестной суммы. В начале второго и третьего годов вносят ещё по $5$ тыс. руб. К концу срока на счёте $118{,}03$ тыс. руб. Определите первоначальную сумму в тысячах рублей.
\item Вклад $300$ тыс. руб. размещён на три года без пополнений и снятий. К концу срока на счёте $399{,}3$ тыс. руб. Найдите годовую процентную ставку.
\end{enumerate}
\subsection*{Сложный уровень}
\begin{enumerate}[resume]
\item Вклад $50$ тыс. руб. размещён на три года под $20\%$ годовых. В начале второго и третьего годов его пополняют на одинаковую сумму. К концу срока накоплено $118{,}08$ тыс. руб. Сколько тысяч рублей вносили каждый раз?
\item Вклад $200$ тыс. руб. размещён на три года под $10\%$ годовых. В начале второго года добавили $20$ тыс. руб., в начале третьего года сняли $10$ тыс. руб. Сколько рублей оказалось на счёте к концу третьего года?
\end{enumerate}

\endgroup
\clearpage
\lessonsection{Ответы к тренировочному блоку}
\renewcommand{\arraystretch}{1.45}
\begin{center}
\begin{tabularx}{0.78\linewidth}{@{}c X@{}}
\toprule
№ & Ответ\\
\midrule
1 & $1100$ руб.\\
2 & $88{,}2$ тыс. руб.\\
3 & $345{,}6$ тыс. руб.\\
4 & $156{,}2$ тыс. руб.\\
5 & $121{,}55$ тыс. руб.\\
6 & $125{,}52$ тыс. руб.\\
7 & $80$ тыс. руб.\\
8 & $10\%$\\
9 & $12$ тыс. руб.\\
10 & $279\,400$ руб.\\
\bottomrule
\end{tabularx}
\end{center}
\vspace{9mm}
\noindent\textcolor{darkaccent}{\textbf{Самопроверка перед сдачей:}}
\begin{itemize}
\item Правильно ли переведена ставка в коэффициент $k$?
\item Учтены ли взносы и снятия \emph{до} начисления процентов соответствующего года?
\item Каждой дополнительной сумме соответствует верное число начислений?
\item Совпадают ли единицы вычисленного результата с единицами, указанными в вопросе?
\end{itemize}
\vfill
\begin{center}
\small\textcolor{darkaccent}{Лёвин Артём Александрович --- эксклюзивно для Ксении Клыковой}
\end{center}
\end{document}
